\section{Power Series}
%%%

If $a_k$ is a sequence of complex numbers, and $z_0\in \CC$, 
the series
\[
 \sum a_k \cdot (z-z_0)^k
\]
dependent on the parameter $z\in \CC$ is called a
\emph{power series}%
\mymargin{power series}%
\index{series!power series}
with coefficients $a_k$ centered at $z_0$.
If we denote by $A\subset \CC$ the set of complex numbers $z$
for which the power series converges
\[
A= \ENCLOSE{z\in \CC \colon \sum a_k (z-z_0)^k \text{ is convergent}}
\]
the sum of the series results in a function
$f \colon A \to \CC$ defined by
\[
  f(z) = \sum_{k=0}^{+\infty} a_k (z-z_0)^k.
\]
The set $A$ is called the \emph{set of convergence}%
\mymargin{set of convergence}%
\index{set!of convergence}
(or \emph{domain of convergence}) of the power series $\sum a_k (z-z_0)^k$.

From now on, to simplify notation, 
we will assume $z_0=0$ so that 
the series becomes $\sum a_k \cdot z^k$. 
In general, one can perform the change of variable $w=z-z_0$ 
to reduce to the case $z_0=0$.

As usual, we may consider power series
with the index $k$ starting from $1$ instead of $0$
(or from any other natural number).
This is not significant, as we can always consider $a_k=0$ for terms
not participating in the summation.

We also observe that for $k=0$ and $z=z_0$, 
the corresponding term of the series is
$a_0 \cdot 0^0 = a_0$ since we have defined $0^0=1$. Thus, we can also write:
\[
  \sum_{k=0}^{+\infty} a_k z^k = a_0 + a_1 z + a_2 z^2 + \dots + a_n z^n + \dots.
\]
Power series thus resemble polynomials, but with infinitely many terms.

\begin{example}[the geometric series]
The power series with coefficients $a_k=1$ is the
geometric series $\sum z^k$.
Thanks to theorem~\ref{th:serie_geometrica} 
we know that the set of convergence is the circle
$A=\ENCLOSE{z\in \CC \colon \abs{z}<1}$
and for $z\in A$ we can calculate the sum of the series:
\[
 f(z) = \sum_{k=0}^{+\infty} z^k  = \frac{1}{1-z}.
\]
If $\abs{z}\ge 1$ the series cannot converge because the term $z^k$ is not
infinitesimal: $\abs{z^k} = \abs{z}^k \ge 1$.
\end{example}

\begin{example}
\label{ex:477474}
The power series $\sum \frac{z^n}{n^n}$ (obtained by setting $a_n=1/n^n$)
has the set of convergence $A=\CC$
since
\[
\sqrt[n]{\abs{\frac{z^n}{n^n}}} = \frac{\abs{z}}{n} \to 0 < 1
\]
and thus by the root test, the series in question converges absolutely
for any $z\in \CC$.
\end{example}

\begin{theorem}[convergence of power series]
\mymark{***}%
If the power series $\sum a_k z^k$ converges at a point $z\in \CC$
(in fact, it suffices that the sequence $a_k z^k$ is bounded)
then the series
converges absolutely for every $w\in \CC$ such that $\abs{w}< \abs{z}$.
Conversely, if the series does not converge at a point $z\in \CC$
then
it does not converge at any $w$ such that $\abs{w} > \abs{z}$ (in fact, the
sequence
$a_n w^n$ is not even bounded, let alone infinitesimal).
\end{theorem}
%
\begin{proof}
\mymark{*}
If the series $\sum a_k z^k$ converges, it means that the sequence
$a_k z^k$ is infinitesimal and, in particular, is bounded.
Thus, there exists $M$ such that for every $k\in \NN$
\[
 \abs{a_k z^k} \le M.
\]
If $z=0$, there is nothing to prove.
If $z\neq 0$, we have
\[
 \abs{a_k} \le \frac{M}{\abs{z}^k}.
\]
Now, chosen any $w\in \CC$ with $\abs{w} < \abs{z}$, we have
\[
  \abs{a_k w^k} = \abs{a_k}\cdot \abs{w}^k \le M \frac{\abs{w}^k}{\abs{z}^k}
  = M q^k
\]
having set $q = \frac{\abs{w}}{\abs{z}}$.
Since $q<1$, the geometric series $\sum q^k$ converges and, by comparison,
the series $\sum \abs{a_k w^k}$ also converges.
Thus, the series $\sum a_k w^k$ converges absolutely.

Conversely, suppose that $\sum a_k z^k$ does not converge
and take $w$ with $\abs{w} > \abs{z}$.
Then $\sum a_k w^k$ cannot converge,
in fact, $a_k w^k$ cannot even be bounded because
if it were, then, by exchanging the roles of $z$ and $w$,
from the previous point, the series $\sum a_k z^k$ would have to converge.
\end{proof}

\begin{corollary}[the set of convergence is circular]%
\mymark{**}
\label{cor:insieme_convergenza}%
Let $\sum a_k z^k$ be a power series and let $A$ be its set of convergence.
Then $A$ is not empty and setting
\[
  R= \sup \ENCLOSE{\abs{z}\colon z\in A}.
\]
it follows that $R\in[0,+\infty]$ and $A$ coincides with the circle centered at
$0$
and of radius $R$ except for the boundary points, in the sense that:
\begin{equation}\label{eq:48463}
   \ENCLOSE{z\in \CC \colon \abs{z} < R}
   \subset A
   \subset \ENCLOSE{z\in \CC \colon \abs{z}\le R}.
\end{equation}
Moreover, the series converges absolutely for every $z\in \CC$ with $\abs{z}<R$
while
the generic term $a_k z^k$ is not even bounded (and hence the series does not
converge)
when $\abs{z}>R$.
\end{corollary}
%
\begin{proof}
If $\abs{z}<R$, it means there exists $z_0\in A$ such that $\abs{z_0} >
\abs{z}$.
But since the series converges at $z_0$ (by definition of $A$), thanks to the
previous theorem, we can assert that the series converges, indeed, converges
absolutely
at $z$. Thus, we obtain the first inclusion in~\eqref{eq:48463}.

If instead we take $z\in \CC$ with $\abs{z}>R$, then
certainly the series does not converge at $z$, otherwise (by the definition of
$R$)
we would have $R\ge \abs{z}$.
Moreover, we can certainly consider a point $z_1\in\CC$
such that $R < \abs{z_1} \le \abs{z}$ and the power series cannot converge at
$z_1$ either.
But then, by the previous theorem,
we deduce that $a_k z^k$ is not even bounded.

We have thus shown the existence of $R\in \bar \RR$ that satisfies
\eqref{eq:48463}.
Clearly, it results $R\ge 0$ because $A$ is a non-empty set
(for every power series, we have $0\in A$).
\end{proof}

\begin{definition}[radius of convergence]
\mymark{**}
The \emph{radius of convergence}%
\mymargin{radius of convergence}%
\index{radius!of convergence} of a power series $\sum a_k z^k$
is the value $R\in[0,+\infty]$
given by corollary~\ref{cor:insieme_convergenza}:
\begin{align*}
    R &= \sup \ENCLOSE{\abs{z}\colon z\in \CC,\ \sum a_k z^k \text{ is
  convergent}}.
\end{align*}
\end{definition}

\begin{theorem}[calculation of the radius of convergence]
\mymark{***}
\label{th:calcolo_raggio_convergenza}
Let $\sum a_n z^n$ be a power series. If the limit exists
\begin{equation}\label{eq:9267345623}
  \lim_{n\to+\infty}\sqrt[n]{\abs{a_n}} =\ell
\end{equation}
then $R=1/\ell$ is the radius of convergence of the series
(where $R=+\infty$ if $\ell = 0$ and $R=0$ if $\ell=+\infty$).

The same holds if the limit exists
\[
  \lim_{n\to +\infty} \frac{\abs{a_{n+1}}}{\abs{a_n}} = \ell.
\]

More generally, it results $R=1/\ell$ if we set
\[
   \ell = \limsup_{n\to +\infty} \sqrt[n]{\abs{a_n}}.
\]
even in the case where the limit in~\eqref{eq:9267345623}
does not exist.
\end{theorem}
%
\begin{proof}
\mymark{***}
Let $r\ge 0$.
Applying the root test to the series $\sum \abs{a_n} r^n$, we have
\[
  \sqrt[n]{\abs{a_n} r^n}
  = r \sqrt[n]{\abs{a_n}} \to r\ell.
\]
Thus, if we choose an $r < 1/\ell$, we have $r \ell<1$ and the series
$\sum a_n z^n$ converges
absolutely for $z=r$.
Therefore, for every $r<1/\ell$,
we find that $r\in A$: it follows that $R\ge 1/\ell$.
If instead we choose $r > 1/\ell$, we have $r \ell > 1$ and thus
$\abs{a_n} r^n \to +\infty$ and the series cannot be convergent
at $z=r$. This means that $R\le 1/\ell$.

The root test also applies in the case where $\ell$ is defined
via $\limsup$.

In the case where the limit of the ratio $\abs{a_{n+1}} / \abs{a_n}$
exists, we know (thanks to the Cesàro convergence criterion
theorem~\ref{th:criterio_cesaro}) that
the limit of the root coincides with the limit of the ratio, and thus
we reduce to the previous case (or we can repeat the proof
using the ratio test instead of the root test).
\end{proof}

\begin{example}
In example~\ref{ex:serie_log}, we have seen
that the set of convergence of the power series
\[
  \sum_{k=1}^{+\infty} \frac{z^k}{k}
\]
is
\[
  A = \ENCLOSE{ z\in \CC \colon \abs{z}\le 1, z\neq 1}.
\]
The radius of convergence must therefore be $R=1$, and this can be
easily verified with one of the previous criteria. For example:
\[
  \lim_{n\to+\infty} \frac{\frac{1}{n+1}}{\frac{1}{n}}
  = \lim_{n\to+\infty} \frac{n}{n+1} = 1
\]
from which $R=1/1=1$.
Note that neither of the two inclusions in \eqref{eq:48463}
is, in this case, an equality.
\end{example}

\begin{theorem}[stability of the radius of convergence]
\label{th:raggio_serie_derivate}
The power series $\sum a_k z^k$ and $\sum k a_k z^k$ have
the same radius of convergence.
\end{theorem}
%
\begin{proof}
  Let $R$ be the radius of convergence of the series $\sum a_k z^k$ 
  and $r$ the radius of convergence of the series $\sum k a_k z^k$.
  
  Since for $k>0$ we have $\abs{a_k z^k} \le \abs{k a_k z^k}$,
  the series $\sum a_k z^k$ converges absolutely at points 
  where the series $\sum k a_k z^k$ converges absolutely.
  This means that $R\ge r$.

  Now, taken a point $w$ with $\abs{w}<R$, consider any
  point $z$ with $\abs{w}<\abs{z}<R$.
  The series $\sum a_k z^k$ converges absolutely at $z$, 
  so $\abs{a_k z^k}$ is infinitesimal and bounded.
  There exists $M$ such that $\abs{a_k z^k}\le M$, 
  which means $\abs{a_k} \le \frac{M}{\abs{z}^k}$.
  But then 
  \[
    \abs{k a_k w^k} 
    \le k \abs{a_k} \abs{w}^k
    \le k M \frac{\abs{w}^k}{\abs{z}^k}
  \]
  from which, setting $q=\frac{\abs w}{\abs z}$, we obtain 
  $\abs{k a_k w^k} \le k M q^k$. 
  Since $q<1$, we know that the series $\sum k M q^k$ 
  is convergent, and thus, by comparison, also the 
  series $\sum k a_k w^k$ is absolutely convergent.
  Since this is true for every $w$ with $\abs{w}<R$,
  it means that $r\ge R$, and this concludes the proof.
\end{proof}
%
\begin{proof}[Alternative Proof]
Since $\sqrt[k]{k}\to 1$, we have
\[
  \limsup \sqrt[k]{\abs{k a_k}} = \limsup \sqrt[k]{\abs{a_k}}.
\]
By theorem~\ref{th:calcolo_raggio_convergenza}, we obtain
that the two corresponding power series have the same radius of convergence.
\end{proof}

\begin{theorem}(continuity of power series)
\label{th:continuita_somma_serie}%
\mymargin{continuity of power series}%
\index{continuity!of power series}%
Let $\sum a_k z^k$ be a power series with radius of convergence $R$.
Then, setting $B=\ENCLOSE{z\in \CC\colon \abs{z}<R}$, the function $f\colon B
\to \CC$
defined by
\[
 f(z) = \sum_{k=0}^{+\infty} a_k z^k
\]
is continuous (on $B$).
\end{theorem}
%
\begin{proof}
Suppose $R>0$, otherwise there is nothing to prove.
Fix $r<R$ and consider two points $z,w\in \CC$ 
with $\abs{z}<r$ and $\abs{w}<r$.
Recall that we have 
\[
  z^n - w^n = (z-w)\cdot(z^{n-1} + wz^{n-2} + \dots + w^{n-2}z + w^{n-1}).
\]
Observing that for each addend on the right side, we have 
$\abs{w^k z^{n-k-1}}\le r^{n-1}$, and since there are $n$ addends, we obtain 
\[
  \abs{z^n - w^n} \le \abs{z-w}\cdot n\cdot  r^{n-1}.
\]
Thus, 
\begin{align*}
  \abs{f(z)-f(w)} 
  &= \abs{\sum_{n=0}^{+\infty} a_n z^n - \sum_{n=0}^{+\infty} a_n w^n}
   = \abs{\sum_{n=0}^{+\infty} a_n (z^n - w^n)} \\
& \le \sum_{n=0}^{+\infty} \abs{a_n} \cdot \abs{z^n - w^n}
\le \abs{z-w}\sum_{n=0}^{+\infty} \abs{a_n} n r^{n-1}.
\end{align*}
But we know that the series $\sum n a_n z^n$ 
has the same radius of convergence $R$ as the original series 
(by theorem~\ref{th:raggio_serie_derivate}),
thus the series $\sum n a_n r^n$ is absolutely convergent and, 
dividing by $r$, 
we find that the series $\sum n a_n r^{n-1}$
is convergent. 
Setting $L = \sum_{n=0}^{+\infty} n \abs{a_n} r^{n-1}$,
from what is written above, we can conclude that 
\[
\abs{f(z)-f(w)} \le L\cdot \abs{z-w}.
\]
Thus, for every $\eps>0$, chosen $\delta = \frac{\eps}{L}$,
if $\abs{z-w}<\delta$, it results $\abs{f(z)-f(w)}\le \eps$.
This means that the function $f$ is continuous at the point $z$, and 
this is true for every $z$ with $\abs{z}<R$.
\end{proof}

The previous theorem guarantees that the sum of a power series
is a continuous function within the radius of convergence.
At points exactly on the boundary of the radius of convergence,
the function $f$ might not be continuous.
However, if the series converges at a
boundary point, the sum of the series is continuous if we approach the point of
convergence along the radius of the disk of convergence, as stated in the
following theorem.

\begin{theorem}[Abel's lemma]
\mymark{*}%
\label{th:lemma_abel}%
\index{lemma!of Abel}%
\index{Abel!lemma of}%
Let
\[
  f(z) = \sum_{k=0}^{+\infty} a_k z^k
\]
be the sum of a power series. 
If the series converges at a point $z_0\in \CC$, $z_0\neq 0$, 
then the series converges for every $z=t z_0$ with $t\in [0,1]$, and moreover,
the function
\[
  t \mapsto f(tz_0)
\]
is continuous at the point $t=1$.
\end{theorem}
%
\begin{proof}
Without loss of generality, we can assume that $f(z_0)=0$; indeed, it suffices
to replace the first term of the series, $a_0$, with $a_0' = a_0 - f(z_0)$ and
prove the theorem for the modified series.
Thus, setting
\[
  A_n = \sum_{k=0}^{n-1} a_k z_0^k
\]
we have that $A_n \to f(z_0) = 0$ and, by definition, $A_0 = 0$.
Using the formula \eqref{eq:somma_per_parti} for summation by parts, for $t\in
[0,1)$
we have
\begin{align*}
\sum_{k=0}^n a_k \cdot (tz_0)^k
= \sum_{k=0}^n a_k z_0^k \cdot t^k
= A_{n+1} \cdot t^{n+1} + \sum_{k=0}^n A_{k+1}\cdot (t^k - t^{k+1})
\end{align*}
and for $n\to +\infty$, we obtain
\[
  f(t\cdot z_0) = f(z_0)\cdot 0 + \sum_{k=0}^{+\infty}A_{k+1}(t^k-t^{k+1})
  = (1-t)\sum_{k=0}^{+\infty} A_{k+1} t^k.
\]
Since $A_n \to 0$, for every $\eps>0$, there exists $m$ such that for every $k >
m$, we have $\abs{A_k} \le  \eps$. Thus, for every $t\in [0,1)$, we have
\begin{align*}
\abs{f(t\cdot z_0)}
 &\le (1-t)\sum_{k=0}^{m-1} \abs{A_{k+1}} t^k
  + (1-t)\sum_{k=m}^{+\infty} \abs{A_{k+1}} t^k \\
  &\le (1-t)\cdot \sum_{k=0}^{m-1}\abs{A_{k+1}} + (1-t)\cdot \eps \cdot
 \frac{t^{m}}{1-t} \\
 &\le (1-t)\sum_{k=0}^{m-1}\abs{A_{k+1}} + \eps.
\end{align*}
Chosen $\delta \le \frac{\eps} {\sum_{k=0}^{m-1} \abs{A_{k+1}}}$, if
$\abs{1-t}<\delta$, we have
\[
  \abs{f(t\cdot z_0)-f(z_0)} = \abs{f(t\cdot z_0)} < 2\eps
\]
which means that $t\mapsto f(t\cdot z_0)$ is continuous at the point $t=1$.
\end{proof}